\VerdierPhysicalPageAnchor{0040}%
\VerdierSemanticAnchor{ch1-1-3-2-eq4}
and the isomorphism $c(g',g)^{-1}$ defines a map

\begin{equation}
  \operatorname{Hom}_{\mathcal C}
    \bigl(T(g)Y,T(g)T(g')Z\bigr)
  \longrightarrow
  \operatorname{Hom}_{\mathcal C}
    \bigl(T(g)Y,T(g'g)Z\bigr).
  \tag{1.3.2.4}
\end{equation}

\VerdierSemanticAnchor{ch1-1-3-2-eq5}
Thus the image of $v$ is an element
$v'\in\operatorname{Hom}_{\mathcal C}(T(g)Y,T(g'g)Z)$. The composite in
$\mathcal C$ of $u$ and $v'$ is an element

\begin{equation}
  w=v'\mathbin{\circ_{\mathcal C}}u,
  \qquad w\in\operatorname{Hom}^{g'g}(X,Z),
  \tag{1.3.2.5}
\end{equation}
which is, by definition, the composite of $u$ and $v$ in
$\mathcal C[G]$. The fibering functor $\mathcal C[G]\to\mathcal G$ is
then the evident functor. The fibered category
$\mathcal C[G]\to\mathcal G$ is equipped with a canonical normalized
cleavage.

\VerdierSemanticAnchor{ch1-1-3-3}
\noindent\textbf{1.3.3.}\quad
There is in fact a more precise result. Let $\mathcal C$ and
$\mathcal C'$ be two $G$-graded categories equipped with normalized
cleavages. The functor $\Psi$ defines a bijection

\VerdierSemanticAnchor{ch1-1-3-3-eq1}
\begin{equation}
  \Psi(\mathcal C,\mathcal C')\colon
  \operatorname{Hom}_{\operatorname{Grad}(G)}(\mathcal C,\mathcal C')
  \longrightarrow
  \operatorname{Hom}_{G^\circ\text{-}\operatorname{Cat}}
    \bigl(\Psi(\mathcal C),\Psi(\mathcal C')\bigr).
  \tag{1.3.3.1}
\end{equation}

These sets are the sets of objects of the categories

\VerdierSemanticAnchor{ch1-1-3-3-eq2}
\begin{equation}
  \mathcal{H}\!om_{\operatorname{Grad}(G)}(\mathcal C,\mathcal C'),
  \qquad
  \mathcal{H}\!om_{G^\circ\text{-}\operatorname{Cat}}
    \bigl(\Psi(\mathcal C),\Psi(\mathcal C')\bigr),
  \tag{1.3.3.2}
\end{equation}
whose morphisms are, respectively, morphisms of functors of degree $e$
and morphisms of functors compatible with the actions of $G^\circ$.
Moreover, it follows immediately from (1.2.7) that the map
$\Psi(\mathcal C,\mathcal C')$ extends to a functor

\VerdierSemanticAnchor{ch1-1-3-3-eq3}
\begin{equation}
  \Psi(\mathcal C,\mathcal C')\colon
  \mathcal{H}\!om_{\operatorname{Grad}(G)}(\mathcal C,\mathcal C')
  \longrightarrow
  \mathcal{H}\!om_{G^\circ\text{-}\operatorname{Cat}}
    \bigl(\Psi(\mathcal C),\Psi(\mathcal C')\bigr)
  \tag{1.3.3.3}
\end{equation}
and one readily checks the following.

\VerdierSemanticAnchor{ch1-1-3-4}
\noindent\textbf{Proposition 1.3.4.}\quad
\emph{The functor $\Psi(\mathcal C,\mathcal C')$ is an isomorphism of
categories.}

\VerdierSemanticAnchor{ch1-1-3-5}
\noindent\textbf{1.3.5.}\quad
Finally, for completeness, one should note that the functor $\Psi$
(extended to morphisms of functors) preserves the convolutions of
functors by morphisms of functors. The verification is left to the
reader.

\VerdierSemanticAnchor{ch1-1-3-6}
\noindent\textbf{1.3.6.}\quad
Most often, we shall use strict $G^\circ$-categories (1.2.2). Let
$(\mathcal C,G^\circ,T)$ and $(\mathcal C',G^\circ,T')$ be two strict
$G^\circ$-categories. A $G^\circ$-functor $(F,m)$ is then a functor
$F\colon\mathcal C\to\mathcal C'$ equipped with isomorphisms

\VerdierSemanticAnchor{ch1-1-3-6-eq1}
\begin{equation}
  m(g)\colon FT(g)\longrightarrow T'(g)F,
  \qquad g\in G,
  \qquad m(e)=\text{identity}
  \tag{1.3.6.1}
\end{equation}
